\begin{proof}

\begin{enumerate}[label={$\langle1\rangle$\arabic*.},leftmargin=*]
\item Push the integral divisor of P4R.6 to
the minimal ruled surface $T_0\to B$.
Its horizontal degree is two. A reduced
horizontal component $H$ satisfies
\[
 0\le 2g(\widetilde H)-2
 \le 2p_a(H)-2
 =-(D_0-H)H\le0.
\]
The first inequality is
Riemann--Hurwitz over the elliptic
base. All inequalities are equalities:
every such component is smooth elliptic,
unramified over $B$, and disjoint from
the other components of $D_0$. Every
vertical fibre meets the horizontal
part, so no vertical component occurs.
The reduced-horizontal possibilities
are therefore one smooth elliptic
bisection or two disjoint elliptic
sections, each with coefficient one.

\item The reverse blowup formula $\mu-1$
then keeps $A=\operatorname{Supp}D$
a disjoint union of their smooth
strict transforms; each new exceptional
curve is absent from $A$. Every
exceptional curve of the normalization
resolution must belong to $A$, by
P4R.6. An exceptional component would
therefore be one of these disjoint
elliptic curves and would be disjoint
from the strict conductor, contrary
to P4R.6. Thus the normalization is
smooth. P4.22, with its hypotheses now
verified, proves minimality and
$e(S)=0$.

\item In the remaining case
$D_0=2H+\sum a_jF_j$ with $H$ a
smooth section. The same reverse
blowups make $A$ the elliptic root $H$
with rational trees attached. Its
inclusion in $T$ induces an
isomorphism on fundamental groups,
by the ruling and invariance under
point blowups. The resolution pushout
gives $\pi_1(S)=\pi_1(C)$ as in
P4R.7. The virtually abelian curve
group calculation gives $e(C)\ge0$,
and P4.13 gives $e(S_{\rm reg})\ge0$.
Their sum proves $e(S)\ge0$ also
in this remaining case.

\end{enumerate}
\end{proof}
